%%%%%%%%%%%%%%%%%%%%%%%%%%%%%%%%%%%%%%%%%%%%%%%%%%%%%%%%%%%%%%%%%%%%%%%%%%%%%%%%
%% 
\cleardoubleoddpage%  Make sure to start each chapter on a new odd page
\chapter{Einleitung}
\label{sec:introduction}


\section{Einleitung}
\label{sec:introduction:motivation}

Die Erstellung von Dienstplänen in Krankenhäusern stellt eine der komplexesten organisatorischen Aufgaben im Gesundheitswesen dar. In den Abteilungen müssen zahlreiche Anforderungen gleichzeitig berücksichtigt werden: Qualifikationen und Verfügbarkeiten des Pflegepersonals, gesetzliche Arbeitszeitregelungen, sowie individuelle Präferenzen der Mitarbeitenden. Kurzfristige Personalausfälle erhöhen die Komplexität zusätzlich und machen eine rein manuelle Planung zunehmend ineffizient. Die händische Erstellung von Dienstplänen ist nicht nur zeitaufwendig, sondern auch fehleranfällig und bietet wenig Skalierbarkeit bei wachsenden Anforderungen.
Die Informatik stellt seit Jahrzehnten formale Methoden zur Lösung derartiger kombinatorischer Optimierungsprobleme bereit. Das sogenannte Nurse Scheduling Problem (NSP) ist in der wissenschaftlichen Literatur gut etabliert und wurde mit einer Vielzahl von Ansätzen untersucht, darunter ganzzahlige lineare Programmierung, metaheuristische Verfahren wie genetische Algorithmen sowie Constraint Programming. In der praktischen Anwendung mangelt es jedoch häufig an Lösungen, die gleichermaßen flexibel gegenüber abteilungsspezifischen Anforderungen, transparent in ihrer Modellierung und nachhaltig wartbar sind.
Answer Set Programming (ASP) bietet in diesem Kontext einen vielversprechenden Ansatz. Als deklaratives Paradigma der Logikprogrammierung ermöglicht ASP eine präzise und gut lesbare Modellierung von Constraints und Optimierungszielen, ohne dass problemspezifische Suchalgorithmen explizit implementiert werden müssen. Mit Clingo, einem der leistungsfähigsten und in der Forschung weitverbreitetsten ASP-Solver, steht ein effizientes Werkzeug zur Verfügung.



\section{Ziel}
\label{sec:introduction:objectives}
Ziel dieser Bachelorarbeit ist die Entwicklung eines ASP-basierten Ansatzes zur automatisierten Erstellung und Optimierung von Dienstplänen in der Augenabteilung eines Krankenhauses, sowie dessen Umsetzung als vollständige Fullstack-Webanwendung. Die Applikation umfasst ein benutzerorientiertes Frontend zur Verwaltung von Personal, Stationen und Kompetenzen sowie ein Backend, das die Optimierungslogik mittels Clingo ausführt und die generierten Dienstpläne strukturiert zurückliefert. Sowohl harte Constraints, wie Mindestbesetzung und Qualifikationsanforderungen, als auch weiche Constraints, wie eine gleichmäßige Verteilung der Zuteilungen zu Stationen, werden dabei formal modelliert und automatisiert gelöst. Die Arbeit untersucht damit, inwieweit ASP als Grundlage für eine praxistaugliche, vollständig integrierte Planungssoftware im klinischen Umfeld eingesetzt werden kann.





%% 
%%%%%%%%%%%%%%%%%%%%%%%%%%%%%%%%%%%%%%%%%%%%%%%%%%%%%%%%%%%%%%%%%%%%%%%%%%%%%%%%
